\documentclass[11pt]{amsart}
\usepackage{amssymb}
\usepackage[margin=1.15in]{geometry}
\usepackage[hidelinks]{hyperref}
\hypersetup{pdftitle={The lazy evaluator of Henderson and Morris, by intersection types}}

% Stamp on every page: file name and compilation date.
\newcommand{\stamptime}{\the\numexpr(\time-30)/60\relax:\ifnum\numexpr\time-60*((\time-30)/60)\relax<10 0\fi\the\numexpr\time-60*((\time-30)/60)\relax}
\AddToHook{shipout/foreground}{%
 \put(0.5\paperwidth,-\paperheight+0.55cm){\makebox[0pt]{\tiny\ttfamily
  unit/doc/\jobname.tex \ --- \ compiled \today, \stamptime}}}

\newtheorem{theorem}{Theorem}[section]
\newtheorem{proposition}[theorem]{Proposition}
\newtheorem{lemma}[theorem]{Lemma}
\newtheorem{corollary}[theorem]{Corollary}
\theoremstyle{definition}
\newtheorem{definition}[theorem]{Definition}
\theoremstyle{remark}
\newtheorem{remark}[theorem]{Remark}

\newcommand{\kw}[1]{\mathsf{#1}}
\newcommand{\at}[1]{\ulcorner #1\urcorner}
\newcommand{\Atom}{\mathsf{atom}}
\newcommand{\obs}[4]{#1\vdash #2:_{#3}#4}
\newcommand{\T}{\mathsf{T}}
\newcommand{\NIL}{\mathsf{NIL}}
\newcommand{\D}{\mathcal D}
\newcommand{\V}{\mathcal V}
\newcommand{\U}{\mathcal U}
\newcommand{\up}{{\uparrow}}

\title[The lazy evaluator by intersection types]{The lazy evaluator of Henderson and Morris,\\ by intersection types}
\date{}

\begin{document}
\begin{abstract}
Henderson and Morris \cite{HM76} give a lazy
evaluator for a pure LISP as seven rewrite rules on a memory, and sketch a
proof that the rules preserve a Scott--Strachey semantics, the delicate
case being the rule that creates a circular structure. We present the
semantics by \emph{intersection types}: the compact elements of the value
domain are the types, ordered by a subtyping relation given by inference
rules, and the meaning of a memory cell is the set of types derivable for
it. All the relations are inductive and positive. There is no
approximation index, no computed join and no consistency condition on the
observations in the invariance proofs (the allocated cells are fresh, as in
the paper). The circular rule is proved by one induction on
derivations, and partial correctness of the evaluator follows.
Finally, the infinite elements are the filters of types. The semantic
function of the paper is then the least solution of its own equations, and
the invariance theorem holds in the paper's form, as equalities of semantic
memories, each rule being a short calculation and the circular rule two
applications of fixed-point induction. Everything is formalised in Unit,
with indexed inductive families (Section~8).
\end{abstract}
\maketitle

\section{The language and the memory}

We follow \S II of the paper. Expressions are
\[
\begin{aligned}
e\ ::=\ {}&\xi \mid (\kw{QUOTE}\ a)\mid(\kw{CONS}\ e\ e)\mid(\kw{CAR}\ e)\mid(\kw{CDR}\ e)\mid(\kw{ATOM}\ e)\mid(\kw{EQ}\ e\ e)\\
\mid{}&(\kw{IF}\ e\ e\ e)\mid(e\ e)\mid(\kw{LAMBDA}\ \xi\ e)\mid(\kw{LABEL}\ \xi\ e)\mid(\kw{FUNARG}\ e\ \theta)
\end{aligned}
\]
with $\xi$ a variable, $a$ an atom and $\theta=\xi_1\,e_1\dots\xi_n\,e_n$ an
association list. A \emph{memory} $\mu$ maps addresses $\pi$ to \emph{cells}:
the same forms with addresses in place of subexpressions. Two memories differ
\emph{at} $\pi_0$ if they agree everywhere except at $\pi_0$; $\mu[e/\pi_0]$
is $\mu$ with the cell $e$ at $\pi_0$.

The rules are those of the paper:
$[C]$ ($\kw{CAR}/\kw{CDR}$ of a $\kw{CONS}$), $[A]$, $[E]$, $[I]$,
$[F]$ ($a$: variable lookup in a $\kw{FUNARG}$, the smallest $i$ with
$\xi_i=\xi$; $b$: a quoted atom; $c$: an application under a $\kw{FUNARG}$;
$d$: a primitive under a $\kw{FUNARG}$), $[G]$ (applying a
$\kw{FUNARG}$-closed $\kw{LAMBDA}$) and $[L]$:
\[
\mu(\pi_0)=(\kw{FUNARG}\ \pi_\ell\ \theta),\quad \mu(\pi_\ell)=(\kw{LABEL}\ \xi\ \pi_1)
\ \Longrightarrow\ \mu'=\mu[(\kw{FUNARG}\ \pi_1\ \xi\ \pi_0\ \theta)/\pi_0],
\]
which creates a circular structure. Each rule rewrites one cell; $[Fc]$ and
$[Fd]$ first allocate new cells. An allocated cell is \emph{fresh}: free, and
occurring in no cell of the memory. In the paper this is part of the
consistency of the memory (``no non-free cell points at a free one''), and
it is the only consistency property used below. The evaluator $\kw{Eval}$ of
the paper performs a sequence of such rewrites.

\section{Types}

The value domain of the paper is $V=A+(V\times V)+(V\to V)$, with Reynolds's
separated sum \cite{Rey71}, so that $\bot$, ``some atom'', $(\kw{CONS}\ \bot\ \bot)$ and
$(\kw{LAMBDA}\ x\ \bot)$ are distinct. Its compact elements are described by
the \emph{types}
\[
\sigma,\tau\ ::=\ \omega\mid \Atom\mid \at{a}\mid \sigma\times\tau\mid \sigma\to\tau\mid \sigma\wedge\tau .
\]
Here $\omega$ is no information, $\Atom$ is ``an atom'', $\at a$ the atom
$a$, $\sigma\times\tau$ a pair whose components have types $\sigma$ and
$\tau$, and $\sigma\to\tau$ a function sending $\sigma$ to $\tau$. Thus
$\omega\times\omega$ is ``a pair'' and $\omega\to\omega$ ``a function''. A finite step function
$\{(u_1,v_1),\dots,(u_n,v_n)\}$ is the type
$(u_1\to v_1)\wedge\dots\wedge(u_n\to v_n)$.

\begin{definition}[subtyping]
$\sigma\le\tau$ is the least relation closed under: reflexivity and
transitivity; $\sigma\wedge\tau\le\sigma$, $\sigma\wedge\tau\le\tau$, and
$\rho\le\sigma,\ \rho\le\tau\ \Rightarrow\ \rho\le\sigma\wedge\tau$;
$\sigma\le\omega$; $\at a\le\Atom$;
$\sigma\le\sigma',\ \tau\le\tau'\Rightarrow\sigma\times\tau\le\sigma'\times\tau'$
and $(\sigma\times\tau)\wedge(\sigma'\times\tau')\le(\sigma\wedge\sigma')\times(\tau\wedge\tau')$;
$\sigma'\le\sigma,\ \tau\le\tau'\Rightarrow\sigma\to\tau\le\sigma'\to\tau'$,
$(\sigma\to\tau)\wedge(\sigma\to\tau')\le\sigma\to(\tau\wedge\tau')$, and
$\sigma\to\omega\le\omega\to\omega$.
\end{definition}

The last axiom says that every arrow type gives ``a function'': all the types
$\sigma\to\omega$ are equivalent to $\omega\to\omega$. It replaces the axiom
$\omega\le\omega\to\omega$ of the usual systems, which would identify ``a
function'' with ``no information'' and so contradict the separated sum.

Presenting compact elements as types is the domain interpretation of
Martin-L\"of \cite{ML83}, also the basis of Abramsky's domain theory in
logical form \cite{Abr91}; for the untyped $\lambda$-calculus it is the
filter model \cite{BCD83}. The order is given by inference rules, so the join
of step functions never has to be computed: it is assembled in a derivation
from $\wedge$ and the arrow rules, and the characterisation of the order
becomes a lemma. Write a type as a meet of \emph{components} (types that are
not meets). A \emph{selection} of components is a finite multiset of them,
possibly empty; the meet of the empty selection is $\omega$.

\begin{lemma}[inversion of subtyping]\label{lem:inv}
Let $\rho$ have components $\rho_k$.
\begin{enumerate}
\item $\rho\le\at a$ iff some $\rho_k$ is $\at a$; $\rho\le\Atom$ iff some
$\rho_k$ is $\Atom$ or an $\at a$.
\item $\rho\le\sigma\times\tau$ iff $\rho$ has a pair component and a
selection of pair components $\sigma_i\times\tau_i$ with
$\bigwedge\sigma_i\le\sigma$ and $\bigwedge\tau_i\le\tau$.
\item ($\beta$-soundness) $\rho\le\sigma\to\tau$ iff $\rho$ has an arrow
component and a selection of arrow components $\sigma_i\to\tau_i$ with
$\sigma\le\bigwedge\sigma_i$ and $\bigwedge\tau_i\le\tau$.
\end{enumerate}
\end{lemma}

\begin{proof}[Proof sketch]
The right-hand sides define a relation $\le^{s}$, \emph{syntax directed}:
inductive on the target and positive (the order appears only in premises).
It contains the generating rules and is reflexive. Transitivity of $\le^s$ is
proved by induction on the depth of the middle type, meets not counted: in
the arrow case, $\rho\le^s\pi$ and $\pi\le^s\sigma\to\tau$ give a selection
of arrows $\pi_k\to\pi'_k$ of $\pi$, each inverted to a selection of arrows of
$\rho$; their union (a meet of selections, which is why selections are
multisets) works, using transitivity through the meets of the $\pi_k$ and of
the $\pi'_k$, whose depth is smaller than that of $\pi$ (their size need not
be, a selection may repeat a component). Hence $\le\ =\ \le^s$.
\end{proof}

Types may be inconsistent ($\at a\wedge\at b$ with $a\ne b$, or
$\Atom\wedge(\omega\times\omega)$). The elements of $V$ are the consistent
filters of types; nothing below needs consistency.

\section{The meaning of a memory cell}

An \emph{environment} $\Gamma$ assigns types to finitely many variables;
$\Gamma(\xi)=\omega$ otherwise, and $\Gamma[\xi:\sigma]$ updates it. For
an association list $\theta=\xi_1\pi_1\dots\xi_n\pi_n$ and types
$\vec\sigma$, write $\Gamma[\theta:\vec\sigma]=\Gamma[\xi_n:\sigma_n]\cdots[\xi_1:\sigma_1]$,
so that the smallest $i$ wins, as in rule $[Fa]$.

\begin{definition}[observation]
$\obs\mu\pi\Gamma\sigma$ (``in $\mu$, the cell $\pi$ has type $\sigma$ in
$\Gamma$'') is the least relation closed under the rules
$\obs\mu\pi\Gamma\omega$;
$\obs\mu\pi\Gamma\sigma,\ \obs\mu\pi\Gamma\tau\Rightarrow\obs\mu\pi\Gamma{\sigma\wedge\tau}$;
$\obs\mu\pi\Gamma\sigma,\ \sigma\le\tau\Rightarrow\obs\mu\pi\Gamma\tau$;
and, according to the cell $\mu(\pi)$:
\begin{itemize}
\item $\xi$: \ $\obs\mu\pi\Gamma{\Gamma(\xi)}$;\quad $(\kw{QUOTE}\ a)$: \ $\obs\mu\pi\Gamma{\at a}$;
\item $(\kw{CONS}\ \pi_1\ \pi_2)$: from $\obs\mu{\pi_1}\Gamma\sigma$ and $\obs\mu{\pi_2}\Gamma\tau$, the type $\sigma\times\tau$;
\item $(\kw{CAR}\ \pi')$: from $\obs\mu{\pi'}\Gamma{\sigma\times\tau}$, the type $\sigma$; $\kw{CDR}$ likewise with $\tau$;
\item $(\kw{ATOM}\ \pi')$: from $\Atom$, the type $\at\T$; from $\omega\times\omega$, the type $\at\NIL$;
\item $(\kw{EQ}\ \pi_1\ \pi_2)$: from $\at{a_1}$ and $\at{a_2}$, the type $\at\T$ if $a_1=a_2$ and $\at\NIL$ otherwise;
\item $(\kw{IF}\ \pi_0\ \pi_1\ \pi_2)$: from $\at\T$ at $\pi_0$ and $\sigma$ at $\pi_1$, the type $\sigma$; from $\at\NIL$ at $\pi_0$ and $\sigma$ at $\pi_2$, the type $\sigma$;
\item $(\pi_0\ \pi_1)$: from $\sigma\to\tau$ at $\pi_0$ and $\sigma$ at $\pi_1$, the type $\tau$;
\item $(\kw{LAMBDA}\ \xi\ \pi')$: from $\obs\mu{\pi'}{\Gamma[\xi:\sigma]}\tau$, the type $\sigma\to\tau$;
\item $(\kw{LABEL}\ \xi\ \pi')$: from $\obs\mu{\pi'}{\Gamma[\xi:\sigma]}\tau$ and $\obs\mu\pi\Gamma\sigma$, the type $\tau$;
\item $(\kw{FUNARG}\ \pi'\ \theta)$: from $\obs\mu{\pi_i}\Gamma{\sigma_i}$ for each binding $\xi_i\pi_i$ of $\theta$ and $\obs\mu{\pi'}{\Gamma[\theta:\vec\sigma]}\tau$, the type $\tau$.
\end{itemize}
\end{definition}

This is the paper's $V\mu\pi\rho$ read as an inductive relation: the meaning
of $\pi$ in an environment is the filter of its types (Section~6 makes this
precise). The paper argues with the approximations $V_i$ of $V$ and $Y_i$ of
the fixed point; here the height of derivations plays their role. Being
inductive, the relation is a least fixed point, so the $\kw{LABEL}$ clause
\emph{is} the paper's $Y\{\lambda x.\ldots\}$; and a circular memory simply
has no derivations that would need infinite unfolding. Note that the
$\kw{LABEL}$ rule has a premise about its own address $\pi$.
Environments are lists of types, not of values. The first three rules make
the set of types of a cell a filter. Two basic facts:

\begin{lemma}\label{lem:basic}
\begin{enumerate}
\item (environments) If $\Delta(\xi)\le\Gamma(\xi)$ for all $\xi$, then
$\obs\mu\pi\Gamma\sigma$ implies $\obs\mu\pi\Delta\sigma$.
\item (generation) A filter containing every type given by one rule of the
cell $\mu(\pi)$ in $\Gamma$, from observations of its components, contains
every type of $\pi$ in $\Gamma$.
\end{enumerate}
\end{lemma}

Both are proved by induction on derivations. Each inversion below is (2)
with a filter chosen for the purpose, and Lemma~\ref{lem:inv}. For example,
a cell $(\kw{CONS}\ \pi_1\ \pi_2)$ has type $\sigma\times\tau$ iff $\pi_1$ has
type $\sigma$ and $\pi_2$ type $\tau$; it has no type $\at a$ and no arrow
type. For a $\kw{FUNARG}$ cell the filter ``the types derived from one choice
of types for $\theta$'' is closed under meets because two choices have a
meet, again a choice, to which Lemma~\ref{lem:basic}(1) moves both
derivations; so every type of $(\kw{FUNARG}\ \pi'\ \theta)$ comes from a
single choice of types for $\theta$.

\section{The rules preserve the meaning}

\begin{theorem}\label{thm:inv}
If $\mu\Rightarrow\mu'$ by one of the rules, then for every non-free cell
$\pi$ of $\mu$, every $\Gamma$ and every $\sigma$:\quad
$\obs\mu\pi\Gamma\sigma\iff\obs{\mu'}\pi\Gamma\sigma$.
\end{theorem}

\noindent\emph{The common part.} Consider first a rewrite of one cell
$\pi_0$. Define $\mu\Rrightarrow\nu$
at $\pi_0$: every type that one rule of the cell $\mu(\pi_0)$ derives from
$\nu$-observations of its components is a $\nu$-type of $\pi_0$. If
$\mu\Rrightarrow\nu$ at $\pi_0$ and $\mu,\nu$ agree elsewhere, then every
$\mu$-observation is a $\nu$-observation: by induction on the derivation,
copying every step that does not use the cell $\pi_0$ and using the hypothesis
for the steps that do. So each direction of each rule is one local lemma.
An allocation keeps the types of the old cells, both ways: forwards the
fresh cell has no rule in $\mu$; backwards, a rule at an old cell has its
premises at its components, which are old since the fresh cell occurs in
no cell of $\mu$, or (for $\kw{LABEL}$) at the cell itself.

The forward directions of $[A]$, $[E]$, $[I]$, PLUS and $[Fb]$ use
Lemma~\ref{lem:inv} to discard impossible cases: $\at a\le\at b$ gives $a=b$,
$\at a$ is below no pair type, and a $\kw{CONS}$ cell has no atom type.

\noindent\emph{The local lemmas.} We give three; the others are similar.
\begin{itemize}
\item $[C]$: $\mu(\pi_0)=(\kw{CAR}\ \pi_1)$, $\mu(\pi_1)=(\kw{CONS}\ \pi_2\ \pi_3)$,
$\mu'=\mu[\mu(\pi_2)/\pi_0]$. A type $\sigma$ of $\pi_0$ by the $\kw{CAR}$ rule
comes from a type $\sigma\times\tau$ of $\pi_1$, hence (generation) from a type
$\sigma$ of $\pi_2$, whose last rule is the rule of the cell $\mu(\pi_2)=\mu'(\pi_0)$.
Conversely, $\pi_0$ and $\pi_2$ hold the same cell in $\mu'$, so they have
the same types, by induction on derivations (a $\kw{LABEL}$ premise at
$\pi_0$ is read at $\pi_2$); and $\pi_3$ has type $\omega$, so $\pi_0$ has type
$\sigma$ by $\kw{CAR}$ of $\sigma\times\omega$. The other copying rules ($[I]$,
$[Fa]$) are treated alike.
\item $[G]$: $\mu(\pi_0)=(\pi_f\ \pi_2)$, $\mu(\pi_f)=(\kw{FUNARG}\ \pi_\lambda\ \theta)$,
$\mu(\pi_\lambda)=(\kw{LAMBDA}\ \xi\ \pi_1)$, $\mu'=\mu[(\kw{FUNARG}\ \pi_1\ \xi\ \pi_2\ \theta)/\pi_0]$.
A type $\tau$ of $\pi_0$ by application comes from $\sigma\to\tau$ at $\pi_f$ and
$\sigma$ at $\pi_2$; by generation, from a single choice $\vec\sigma$ of types
for $\theta$ and, by Lemma~\ref{lem:inv}(3), from $\pi_1:\tau$ in
$\Gamma[\theta:\vec\sigma][\xi:\sigma]$, the environment of the new
$\kw{FUNARG}$ cell with the choice $\vec\sigma$ updated at $\xi$. Conversely,
a derivation at the new cell gives $\pi_2:\sigma$ and a choice for $\theta$;
the choice for the old $\kw{FUNARG}$ cell must give $\omega$ to $\xi$, since
its binding of $\xi$ is shadowed by the new one.
\item $[Fa]$: $\mu(\pi_0)=(\kw{FUNARG}\ \pi'\ \theta)$, $\mu(\pi')=\xi$, $i$ the
smallest index with $\xi_i=\xi$, $\mu'=\mu[\mu(\pi_i)/\pi_0]$. A type of $\pi_0$
is above $\Gamma[\theta:\vec\sigma](\xi)=\sigma_i$, a type of $\pi_i$; and
conversely.
\item $[Fc]$, $[Fd]$: after the allocation of fresh cells $\pi'_j$ holding
$(\kw{FUNARG}\ \pi_j\ \theta)$, the cell $(\kw{FUNARG}\ (\pi_1\ \pi_2)\ \theta)$
becomes $(\pi'_1\ \pi'_2)$. Forwards, the types chosen for $\theta$ at $\pi_0$
serve for both new cells. Backwards, the two new cells may choose different
types for $\theta$; their meet is a choice of types for $\theta$, and
Lemma~\ref{lem:basic}(1) moves both subderivations to it.
\end{itemize}

\noindent\emph{The circular rule $[L]$.} Here
$\mu(\pi_0)=(\kw{FUNARG}\ \pi_\ell\ \theta)$, $\mu(\pi_\ell)=(\kw{LABEL}\ \xi\ \pi_1)$
and $\nu=\mu[(\kw{FUNARG}\ \pi_1\ \xi\ \pi_0\ \theta)/\pi_0]$.

($\mu\Rightarrow\nu$.) By induction on the derivation of
$\obs\mu{\pi_\ell}{\Gamma[\theta:\vec\sigma]}\tau$, with $\vec\sigma$ types of
$\theta$ in $\nu$, using the induction hypothesis on the self-premise of the
$\kw{LABEL}$ rule: that rule uses
$\pi_1:\tau$ in $\Gamma[\theta:\vec\sigma][\xi:\sigma]$ and a strictly smaller
derivation of $\pi_\ell:\sigma$, which by induction gives $\pi_0:\sigma$ in $\nu$;
these are exactly the premises of the $\kw{FUNARG}$ rule of $\nu(\pi_0)$, the
binding $\xi\pi_0$ being first. Each unfolding of the label becomes one step
through the circular binding: this single induction replaces the paper's two
nested approximations $V_i$ and $Y_i$.

($\nu\Rightarrow\mu$.) By the common part, it suffices that a type $\tau$
given at $\pi_0$ by the $\kw{FUNARG}$ rule of $\nu(\pi_0)$, from
$\mu$-observations, is a $\mu$-type of $\pi_0$. That rule gives a choice
$\vec\sigma$ for $\theta$, a type $\sigma$ of $\pi_0$ for $\xi$, and $\pi_1:\tau$ in
$\Gamma[\theta:\vec\sigma][\xi:\sigma]$. In $\mu$, by generation, $\pi_0:\sigma$
comes from a choice $\vec\sigma'$ for $\theta$ and $\pi_\ell:\sigma$ in
$\Gamma[\theta:\vec\sigma']$. Their meet $\vec\delta$ is a choice for $\theta$, and by
Lemma~\ref{lem:basic}(1) both derivations hold with $\vec\delta$; the $\kw{LABEL}$
rule then gives $\pi_\ell:\tau$ in $\Gamma[\theta:\vec\delta]$, hence $\pi_0:\tau$.
The meet always exists, inconsistent or not: no consistency of the
observations of a cell is needed.

\section{Partial correctness}

\begin{corollary}\label{cor:pc}
Let $\kw{Eval}(\pi_0,\mu)$ halt with $\mu'$, $\pi_0$ non-free in $\mu$. Then
$\pi_0$ has the same types in $\mu$ and $\mu'$, in every environment, and:
if $\mu'(\pi_0)=(\kw{QUOTE}\ a)$, the types of $\pi_0$ are exactly those above
$\at a$: its value is the atom $a$;
if $\mu'(\pi_0)$ is a $\kw{CONS}$, $\pi_0$ has type $\omega\times\omega$ and no
atom or arrow type: its value is a pair;
if $\mu'(\pi_0)=(\kw{FUNARG}\ (\kw{LAMBDA}\dots)\dots)$, it has type
$\omega\to\omega$ and no atom or pair type: its value is a function.
\end{corollary}

\begin{proof}
$\kw{Eval}$ performs a finite sequence of allocations and rule
applications, so Theorem~\ref{thm:inv} applies at each step; $\pi_0$ stays
non-free, since every rewrite writes a non-free cell (the copies are of
evaluated cells). The three cases are Lemma~\ref{lem:basic}(2) at the final
cell, with Lemma~\ref{lem:inv}.
\end{proof}

As in the paper, the same argument shows that the final form does not depend
on the order of the rewrites. No consistency of the memory is used: the
meaning of $\pi_0$ is shown equal to a principal filter of a consistent type.
Example~4 of the paper (\texttt{integers}, evaluating to $1$) is the instance
where $\mu'(\pi_0)=(\kw{QUOTE}\ 1)$.

\section{Infinite elements: the semantics of Henderson and Morris}

The preceding sections use only the compact elements. We now add the
infinite elements, as filters of types, and recover the paper's statements
literally: its semantic function is the least solution of its own
equations, and each rule leaves the semantic memory unchanged. This is the
construction of the elements as filters of compact elements
\cite{ML83,Sco82,Abr91}, as in the filter model \cite{BCD83}, here for the
domain with separated sum.

\subsection{The domain}

A \emph{filter} is a set $x$ of types with $\omega\in x$, closed under
$\wedge$ and upward closed for $\le$. Let $\D$ be the set of filters,
ordered by inclusion. It is a complete lattice: the intersection of filters
is a filter, and so is a directed union. Its least element is
$\bot=\up\omega=\{\sigma\mid\omega\le\sigma\}$; its greatest element is the
set of all types. The \emph{compact} elements are the principal filters
$\up\sigma$, and every filter is the directed union of the $\up\sigma$,
$\sigma\in x$. We keep all the filters, the inconsistent ones included:
$\D$ is a lattice, as the paper's domain is (``$\bot$ being the bottom of
the whole lattice''), and nothing requires consistency. Two filters are
equal when they have the same elements.

For a set $A$ of types, the filter \emph{generated} by $A$ is
$\langle A\rangle=\{\tau\mid\omega\le\tau$, or
$\sigma_1\wedge\dots\wedge\sigma_n\le\tau$ for some
$\sigma_1,\dots,\sigma_n\in A$, $n\ge1\}$.
The operations of the paper become, for $x,y,z\in\D$:
\[
\begin{aligned}
\langle x,y\rangle&=\langle\{\sigma\times\tau\mid\sigma\in x,\ \tau\in y\}\rangle, &
x.1&=\{\sigma\mid\omega\le\sigma\ \text{or}\ \sigma\times\omega\in x\},\quad
x.2\ \text{likewise},\\
x\cdot y&=\{\tau\mid\omega\le\tau\ \text{or}\ \sigma\to\tau\in x\ \text{for some}\ \sigma\in y\}, &
\kw{atom}(x)&=\langle\{\at\T\mid\Atom\in x\}\cup\{\at\NIL\mid\omega\times\omega\in x\}\rangle,
\end{aligned}
\]
$\kw{eq}(x,y)$ is generated by $\at\T$ if $\at a\in x$ and $\at a\in y$ for
some atom $a$, and by $\at\NIL$ if $\at a\in x$ and $\at b\in y$ with
$a\ne b$; $\kw{if}(x,y,z)$ is generated by $y$ if $\at\T\in x$ and by $z$ if
$\at\NIL\in x$. Each is a filter (for $x\cdot y$ by the two arrow rules of
Definition~2.1), and each is monotone in its arguments. The summand
$\omega\le\tau$ is the separated sum: $\bot\cdot y=\bot$, but
$\langle\bot,\bot\rangle=\up(\omega\times\omega)\ne\bot$.

A monotone $g:\D\to\D$ is \emph{continuous} if
$g(x)=\bigcup_{\sigma\in x}g(\up\sigma)$; it is then determined by its
\emph{trace}, the relation $\tau\in g(\up\sigma)$ between types. For such
$g$ put
\[
\Lambda g=\langle\{\sigma\to\tau\mid\tau\in g(\up\sigma)\}\rangle,
\]
and let $Yg$ be the least set of types containing $\omega$, closed under
$\wedge$ and $\le$, and such that $\sigma\in Yg$ and $\tau\in g(\up\sigma)$
give $\tau\in Yg$. Both depend only on the trace of $g$. Note
$\omega\to\omega\in\Lambda g$, so $\Lambda g\ne\bot$: ``$(\kw{LAMBDA}\ x\ \bot)$''
and $\bot$ are distinct, as the paper requires.

\begin{lemma}[laws]\label{lem:laws}
For $x,y,z\in\D$, atoms $a,b$, and $g$ continuous:
\begin{enumerate}
\item $\langle x,y\rangle.1=x$ and $\langle x,y\rangle.2=y$;
\item ($\beta$) $(\Lambda g)\cdot y=g(y)$;
\item $Yg=g(Yg)$, and $g(x)\subseteq x$ implies $Yg\subseteq x$;
\item $\kw{atom}(\up\at a)=\up\at\T$, $\kw{atom}\langle x,y\rangle=\up\at\NIL$;
$\kw{eq}(\up\at a,\up\at a)=\up\at\T$, $\kw{eq}(\up\at a,\up\at b)=\up\at\NIL$
for $a\ne b$; $\kw{if}(\up\at\T,y,z)=y$, $\kw{if}(\up\at\NIL,y,z)=z$;
\item $\up\at a$, $\langle x,y\rangle$ and $\Lambda g$ are pairwise distinct
and distinct from $\bot$: they contain respectively $\at a$, $\omega\times\omega$,
$\omega\to\omega$, and no other of these three types.
\end{enumerate}
\end{lemma}

\begin{proof}
All but (3) are Lemma~\ref{lem:inv} read on generated filters. For (1): if
$\sigma\times\omega\in\langle x,y\rangle$ and $\omega\not\le\sigma$, then
$\bigwedge(\sigma_i\times\tau_i)\le\sigma\times\omega$ with $\sigma_i\in x$,
so $\bigwedge\sigma_i\le\sigma$ and $\sigma\in x$. For (2): let
$\sigma\to\tau\in\Lambda g$ with $\sigma\in y$; then
$\bigwedge_k(\sigma_k\to\tau_k)\le\sigma\to\tau$ with
$\tau_k\in g(\up\sigma_k)$ ($\omega\le\sigma\to\tau$ is impossible), and
Lemma~\ref{lem:inv}(3) gives a selection with $\sigma\le\sigma_k$ and
$\bigwedge\tau_k\le\tau$; since $\up\sigma_k\subseteq\up\sigma\subseteq y$,
monotonicity gives $\tau\in g(y)$.
Conversely $\tau\in g(y)$ is in some $g(\up\sigma)$, $\sigma\in y$, and
$\sigma\to\tau\in\Lambda g$. For (3): if $g(x)\subseteq x$, then $Yg\subseteq x$
by induction on the generation of $Yg$, since $\sigma\in x$ and
$\tau\in g(\up\sigma)$ give $\tau\in g(x)\subseteq x$. Next
$Yg\subseteq g(Yg)$ by the same induction, $\up\sigma\subseteq Yg$ giving
$g(\up\sigma)\subseteq g(Yg)$; and $g(Yg)\subseteq Yg$ by continuity and the
last closure rule.
\end{proof}

The least fixed point is thus obtained by an inductive generation, not as
the limit of the iterates $g^n(\bot)$; (3) is fixed-point induction
(Park's principle \cite{Par69}).

\subsection{Semantic memories}

An \emph{environment} $\rho$ maps variables to filters. A type
environment $\Gamma$ (Section~3) is \emph{in} $\rho$ when
$\Gamma(\xi)\in\rho(\xi)$ for all $\xi$; write $\up\Gamma$ for the
environment $\xi\mapsto\up\Gamma(\xi)$. A function $F$ from environments to
filters is continuous if it is monotone and
$F(\rho)=\bigcup_{\Gamma\in\rho}F(\up\Gamma)$. A \emph{semantic memory} $S$
maps each address to a continuous function from environments to filters.
Following the paper, for a cell $e$ let $\U(S)(e)\rho$ be:
\[
\begin{array}{ll}
\rho(\xi) & e=\xi\\
\up\at a & e=(\kw{QUOTE}\ a)\\
\langle S\pi_1\rho,S\pi_2\rho\rangle & e=(\kw{CONS}\ \pi_1\ \pi_2)\\
(S\pi'\rho).1,\quad(S\pi'\rho).2 & e=(\kw{CAR}\ \pi'),\ (\kw{CDR}\ \pi')\\
\kw{atom}(S\pi'\rho),\ \kw{eq}(S\pi_1\rho,S\pi_2\rho),\ \kw{if}(S\pi_0\rho,S\pi_1\rho,S\pi_2\rho)
 & e=(\kw{ATOM}\ \pi'),\ (\kw{EQ}\ \pi_1\ \pi_2),\ (\kw{IF}\ \pi_0\ \pi_1\ \pi_2)\\
S\pi_0\rho\cdot S\pi_1\rho & e=(\pi_0\ \pi_1)\\
\Lambda(x\mapsto S\pi'\rho[x/\xi]) & e=(\kw{LAMBDA}\ \xi\ \pi')\\
Y(x\mapsto S\pi'\rho[x/\xi]) & e=(\kw{LABEL}\ \xi\ \pi')\\
S\pi'\rho[S\theta\rho/\theta] & e=(\kw{FUNARG}\ \pi'\ \theta)\\
\bot & \pi\ \text{free}
\end{array}
\]
where $\rho[S\theta\rho/\theta]$ binds each $\xi_i$ of $\theta=\xi_1\pi_1\dots\xi_n\pi_n$
to $S\pi_i\rho$, the smallest $i$ winning. Each clause is continuous in
$\rho$ (for $Y$, by induction on its generation). The paper defines its semantic function $\V$ as the least solution
of $\V\mu\pi=\U(\V\mu)(\mu\pi)$. We take the observations of Section~3 as
the definition:
\[
\V\mu\pi\rho=\{\sigma\mid\obs\mu\pi\Gamma\sigma\ \text{for some}\ \Gamma\ \text{in}\ \rho\}.
\]

\begin{theorem}\label{thm:fix}
\begin{enumerate}
\item $\V\mu\pi\rho$ is a filter, continuous in $\rho$, and
$\sigma\in\V\mu\pi(\up\Gamma)$ iff $\obs\mu\pi\Gamma\sigma$.
\item $\V\mu\pi=\U(\V\mu)(\mu\pi)$ for every $\pi$.
\item If $S$ is a semantic memory with $\U(S)(\mu\pi)\subseteq S\pi$ for all
$\pi$, then $\V\mu\pi\subseteq S\pi$ for all $\pi$.
\end{enumerate}
\end{theorem}

\begin{proof}
(1) $\omega$ is a type of $\pi$ in the empty environment. If $\Gamma_1$ and
$\Gamma_2$ are in $\rho$, so is their pointwise meet, and
Lemma~\ref{lem:basic}(1) moves both derivations to it: this gives the meet
rule and continuity. The last claim is Lemma~\ref{lem:basic}(1).

($\supseteq$ in (2)) Every generator of $\U(\V\mu)(\mu\pi)\rho$ is the
conclusion of the rule of the cell $\mu\pi$, in an environment in $\rho$
obtained as the meet of the environments of the premises. For
$\kw{LAMBDA}$: if $\tau\in\V\mu\pi'(\rho[\up\sigma/\xi])$, by $\Gamma'$, then
$\Gamma=\Gamma'[\xi:\omega]$ is in $\rho$ and $\Gamma[\xi:\sigma]$ is below
$\Gamma'$, so $\pi:\sigma\to\tau$ in $\Gamma$. For $\kw{LABEL}$ use
Lemma~\ref{lem:laws}(3) with $x=\V\mu\pi\rho$: by continuity, an element of
$\V\mu\pi'(\rho[x/\xi])$ is in some $\V\mu\pi'(\rho[\up\sigma/\xi])$ with
$\sigma\in x$, and the $\kw{LABEL}$ rule applies. For $\kw{FUNARG}$ the
environment of the premises meets the environments given for the
$\pi_i$; this is where the meet of the choices of types for $\theta$
(Section~4) now happens, inside the filters.

($\subseteq$ in (2), and (3)) Let $S$ satisfy $\U(S)(\mu\pi)\subseteq S\pi$
for all $\pi$. By induction on derivations, $\obs\mu\pi\Gamma\sigma$ implies
$\sigma\in\U(S)(\mu\pi)\rho$ for every $\rho$ containing $\Gamma$: the
premises at a component $\pi_i$ give, by induction, types in
$\U(S)(\mu\pi_i)\subseteq S\pi_i$, and the rule is the corresponding
generator. For $\kw{LABEL}$ the premise $\obs\mu\pi\Gamma\sigma$ at the cell
itself gives, by induction, $\sigma\in\U(S)(\mu\pi)\rho=Y(\dots)$, and with
$\tau\in S\pi'(\rho[\up\sigma/\xi])$ the closure rule of $Y$ gives
$\tau\in Y(\dots)$. Hence $\V\mu\pi\subseteq\U(S)(\mu\pi)\subseteq S\pi$.
Taking $S=\V\mu$, a prefixed point by the first part, gives $\subseteq$
in (2).
\end{proof}

So $\V\mu$ is the least fixed point of $S\mapsto\U(S)(\mu\,\cdot)$, that is,
the paper's $\V\mu=\lim_i\V_i\mu$; the approximations are not needed. By
(1), Theorem~\ref{thm:inv} is the statement $\V\mu\pi=\V\mu'\pi$ restricted to
the environments $\up\Gamma$, and continuity gives the converse. We now prove
the equalities directly, as the paper does.

\subsection{The rules leave the semantic memory unchanged}

\begin{lemma}[comparison]\label{lem:cmp}
Let $\nu=\mu[e/\pi_0]$. If $\U(\V\nu)(\mu\pi_0)\subseteq\V\nu\pi_0$, then
$\V\mu\subseteq\V\nu$.
\end{lemma}

\begin{proof}
For $\pi\ne\pi_0$, $\U(\V\nu)(\mu\pi)=\U(\V\nu)(\nu\pi)=\V\nu\pi$ by
Theorem~\ref{thm:fix}(2); so $\V\nu$ satisfies the hypothesis of
Theorem~\ref{thm:fix}(3) for $\mu$.
\end{proof}

So for each rule two inclusions at $\pi_0$ are to be checked, and for all
rules but $[L]$ both follow from one equation: if
$\U(S)(\mu\pi_0)=\U(S)(\nu\pi_0)$ for every semantic memory $S$ satisfying
$S\pi=\U(S)(\mu\pi)$ at the cells $\pi\ne\pi_0$, then, taking $S=\V\nu$ and
$S=\V\mu$, $\V\mu=\V\nu$. Note that $\U(S)(e)$ does not mention the address
of $e$; this is why the copying rules need no care about $\kw{LABEL}$ cells
here. The equations, with $\rho'=\rho[S\theta\rho/\theta]$:
\begin{itemize}
\item $[C]$: $\U(S)(\kw{CAR}\ \pi_1)\rho=(S\pi_1\rho).1=\langle S\pi_2\rho,S\pi_3\rho\rangle.1
=S\pi_2\rho=\U(S)(\mu\pi_2)\rho$, by Lemma~\ref{lem:laws}(1). (If $\pi_2=\pi_0$
then $\nu=\mu$.)
\item $[A]$, $[E]$, $[I]$: Lemma~\ref{lem:laws}(4), in the same way.
\item $[Fa]$: $\U(S)(\kw{FUNARG}\ \pi'\ \theta)\rho=S\pi'\rho'=\rho'(\xi)=S\pi_i\rho=\U(S)(\mu\pi_i)\rho$;
$[Fb]$: $S\pi'\rho'=\up\at a$.
\item $[G]$: $\U(S)(\pi_f\ \pi_2)\rho=S\pi_f\rho\cdot S\pi_2\rho
=\Lambda(x\mapsto S\pi_1\rho'[x/\xi])\cdot S\pi_2\rho
=S\pi_1\rho'[S\pi_2\rho/\xi]=\U(S)(\kw{FUNARG}\ \pi_1\ \xi\,\pi_2\,\theta)\rho$,
by $\beta$ (Lemma~\ref{lem:laws}(2)).
\end{itemize}

\noindent\emph{The circular rule $[L]$}, with
$\mu\pi_0=(\kw{FUNARG}\ \pi_\ell\ \theta)$, $\mu\pi_\ell=(\kw{LABEL}\ \xi\ \pi_1)$,
$\nu\pi_0=(\kw{FUNARG}\ \pi_1\ \xi\,\pi_0\,\theta)$. Write
$g_S(x)=S\pi_1\rho'[x/\xi]$. Then $\U(S)(\mu\pi_0)\rho=S\pi_\ell\rho'=Yg_S$
and $\U(S)(\nu\pi_0)\rho=g_S(S\pi_0\rho)$.
\begin{itemize}
\item ($\V\mu\subseteq\V\nu$) For $S=\V\nu$: $g_S(\V\nu\pi_0\rho)=\U(\V\nu)(\nu\pi_0)\rho=\V\nu\pi_0\rho$,
so $Yg_S\subseteq\V\nu\pi_0\rho$ by Lemma~\ref{lem:laws}(3).
\item ($\V\nu\subseteq\V\mu$) For $S=\V\mu$: $\V\mu\pi_0\rho=Yg_S$, so
$\U(\V\mu)(\nu\pi_0)\rho=g_S(Yg_S)=Yg_S=\V\mu\pi_0\rho$.
\end{itemize}
The paper proves the second inclusion by cutting the loop of $Y$ and the
loop of $\V$ together (approximations $\V_i$ and $Y_i$); with both
characterised as least fixed points, each inclusion is one application of
fixed-point induction.

\noindent\emph{The allocating rules $[Fc]$, $[Fd]$}: $\mu\pi_0=(\kw{FUNARG}\ \pi_a\ \theta)$
with $\mu\pi_a$ a primitive form $p(\pi_1,\dots,\pi_k)$ (an application or a
primitive). The rule first allocates cells $\pi'_j$ holding
$(\kw{FUNARG}\ \pi_j\ \theta)$, then rewrites $\pi_0$ to $p(\pi'_1,\dots,\pi'_k)$.
The rewrite is a one-cell rule: with $P$ the operation of $p$,
$\U(S)(\mu\pi_0)\rho=S\pi_a\rho'=P(S\pi_1\rho',\dots)=P(S\pi'_1\rho,\dots)$,
since $S\pi'_j\rho=S\pi_j\rho'$. An allocation $\nu=\mu[(\kw{FUNARG}\ \pi_j\ \theta)/\pi']$
takes a free cell $\pi'$ occurring in no cell of $\mu$ (as in the paper, no
non-free cell points at a free one). Then $\V\mu\subseteq\V\nu$ by
Lemma~\ref{lem:cmp}, since $\U(\V\nu)(\mu\pi')=\bot$. Conversely, let $S$ be
$\V\mu$ except at $\pi'$, where it is the filter of all types. No cell of
$\mu$ mentions $\pi'$, so $\U(S)(\nu\pi)=\U(\V\mu)(\mu\pi)\subseteq S\pi$
for $\pi\ne\pi'$, and $\U(S)(\nu\pi')\subseteq S\pi'$ trivially. Hence
$\V\nu\subseteq S$ by Theorem~\ref{thm:fix}(3), and $\V\nu\pi=\V\mu\pi$ at
the non-free cells of $\mu$.

\begin{theorem}\label{thm:hm}
If $\mu\Rightarrow\nu$ by one of the rules, then $\V\mu\pi=\V\nu\pi$ for
every non-free cell $\pi$ of $\mu$. If $\kw{Eval}(\pi_0,\mu)$ halts, then
$\V\mu\pi_0\bot$ is $\up\at a$, a pair $\langle x,y\rangle$ or a function
$\Lambda g$, according as it halts with $(\kw{QUOTE}\ a)$, a $\kw{CONS}$ or a
$\kw{FUNARG}$ of a $\kw{LAMBDA}$ at $\pi_0$. In particular, if
$\V\mu\pi_0\bot=\bot$ then $\kw{Eval}(\pi_0,\mu)$ does not halt, and if
$\V\mu\pi_0\bot=\up\at a$ and it halts, it halts with $(\kw{QUOTE}\ a)$.
\end{theorem}

\noindent These are the paper's four correctness clauses; the last two
sentences are Lemma~\ref{lem:laws}(5), which rests on the separated sum.

\section{Relation to the finite-element presentation}

An earlier formalisation, in Agda, presents the compact elements as finite
elements, ordered by an algorithm. Up to the equivalence $\sigma\le\tau\le\sigma$, the consistent types are the
finite elements: $\omega$ is $\bot$, $\at a$ an atom, $\sigma\times\tau$ a pair
and $\bigwedge(u_i\to v_i)$ the step function $\{(u_i,v_i)\}$. Subtyping is the
\emph{reverse} of the information order: $u\sqsubseteq v$ iff $v\le u$, so
$\omega$, the top type, is the least element $\bot$, and the meet of types is
the join of the corresponding elements, when they are consistent. What changes is
the presentation of the order: by inference rules instead of an algorithm
computing $h(u)$. This removes the stage stratification of the order and its
stability lemmas, the duplicate (propositional and decidable) orders, the
consistency proof of observations, and the ``shape-directed'' inversion forced
by the closure rules (here: Lemma~\ref{lem:basic}(2)). The order algorithm, if
wanted for computation, can be added afterwards as a decision procedure proved
equivalent to $\le$.

\section{Formalisation in Unit}

Types are one ordinary inductive type. Subtyping $\le$, its syntax-directed
form $\le^s$ (with a family of selections of arrow components), and the
observation relation are indexed inductive families \cite{CPM90}, with all
recursive premises positive; none of them is mutual or nested. The rule
instances of a cell are separated from their premises: the premises form a
function field indexed by a type of premises, which gives the
$\kw{FUNARG}$ rule one premise per binding. The memory is a parameter (a
function from addresses to cells), the environment a list of types. The
evaluator is a relation (big-step, as $\kw{Eval}$), each case restating the
cells of its rewrite in the memory reached. Sections~2--5 are formalised in
\texttt{domain/lazy}, with no parameters; for Example~4 the observation
derivation, not the paper's run of the evaluator.

The infinite elements of Section~6 need function extensionality and
propositional extensionality (both consequences of univalence) to state
equalities of filters and of semantic memories, and propositional
truncation for ``for some $\Gamma$ in $\rho$''. They are parameters of the
development, which the finite part above does not use. A filter is a
proposition-valued predicate on types with the three closure properties;
the filters form a type one universe up. Continuous functions are given by
their traces, relations between types, so that $\Lambda$ and $Y$ are
inductive filters in the universe of the types and no resizing is needed.
This part is formalised in \texttt{domain/ideal} (filters and operations,
Lemma~\ref{lem:laws}, Theorem~\ref{thm:fix}, Lemma~\ref{lem:cmp} and
Theorem~\ref{thm:hm}); its proof of Theorem~\ref{thm:hm} uses the
observations only through Theorem~\ref{thm:fix}, not
Theorem~\ref{thm:inv}. In the converse of Theorem~\ref{thm:fix}(1), a
witness $\Gamma'$ in $\up\Gamma$ gives $\Gamma(\xi)\le\Gamma'(\xi)$ for each
$\xi$ only as a proposition; since $\Gamma'$ is a finite list and
$\Gamma'(\xi)=\omega$ at the other variables, these are obtained one binding
at a time, and Lemma~\ref{lem:basic}(1) moves the observation to $\Gamma$.

\begin{thebibliography}{99}
\bibitem{HM76} P. Henderson and J. H. Morris, A lazy evaluator.
\emph{Proc. 3rd ACM Symposium on Principles of Programming Languages}, 1976, 95--103.
\bibitem{ML83} P. Martin-L\"of, The domain interpretation of type theory.
Lecture notes, Workshop on Semantics of Programming Languages (K. Karlsson,
K. Petersson, eds.), Chalmers University of Technology and University of
G\"oteborg, August 1983.
\bibitem{Abr91} S. Abramsky, Domain theory in logical form.
\emph{Annals of Pure and Applied Logic} 51 (1991), 1--77; LICS 1987.
\bibitem{BCD83} H. Barendregt, M. Coppo and M. Dezani-Ciancaglini, A filter
lambda model and the completeness of type assignment.
\emph{Journal of Symbolic Logic} 48 (1983), 931--940.
\bibitem{Sco82} D. Scott, Domains for denotational semantics.
\emph{ICALP 1982}, LNCS 140, 577--613.
\bibitem{Rey71} J. C. Reynolds, Notes on a lattice-theoretic approach to the
theory of computation. Lecture notes, Syracuse University, 1971.
\bibitem{Par69} D. Park, Fixpoint induction and proofs of program properties.
\emph{Machine Intelligence} 5 (1969), 59--78.
\bibitem{CPM90} Th. Coquand and C. Paulin-Mohring, Inductively defined types.
\emph{COLOG-88}, LNCS 417 (1990), 50--66.
\end{thebibliography}

\end{document}
